% !PRINTED 612
\weberpage{612}

%<<CONT
so wird

\[ \chi = \frac{Q\varphi}{\xi} = \beta_1t_1+\beta_2t_2, \]

und folglich sind $\beta_1,\beta_2$ ganze durch $\psi'$ theilbare Zahlen.

Bilden wir aus (9) die Discriminante $\varDelta(\beta_1,\beta_2)$, so erhalten wir

\weq{10}{\varDelta(\beta_1,\beta_2) = \begin{vmatrix}\beta_1, & \beta_2\\ \beta_1', & \beta_2'\end{vmatrix}^2 = \frac{e^4c^4Q^4}{(\xi\xi')^2}(\omega-\omega')^2.}

Aus (5) und (7) folgt aber

\[ \xi\xi' = N(\xi) = PQ = e^2cQ, \]

und wenn wir noch $\omega-\omega' = \dfrac{\sqrt{\varDelta}}{c}$ setzen, so findet sich

\weq{11}{\varDelta(\beta_1,\beta_2) = Q^2\varDelta.}

Hieraus folgt aber nach dem Satze \S.\,147, 2., dass $\beta_1,\beta_2$ eine Basis von $\psi'$ ist.

Andererseits ist nach (6) auch $e'c', e'c'\eta'$ eine Basis von $\psi'$, und folglich lassen sich die ganzen rationalen Zahlen $p,q,r,s$ so bestimmen, dass

\weq{12}{\begin{aligned}
\beta_1 &= \frac{Qec}{\xi} = e'c'(r\eta'+s)\\
\beta_2 &= \frac{Qec}{\xi}\,\omega = e'c'(p\eta'+q),
\end{aligned}}

und zugleich

\weq{13}{ps-qr = \pm 1.}

Aus (12) folgt aber durch Division:

\weq{14}{\omega = \frac{p\eta'+q}{r\eta'+s},}

d. h. $\omega$ ist mit $\eta'$ äquivalent.

Ersetzen wir nun $\varphi$ durch die äquivalente Form $\varphi_1$, so besteht eine Gleichung wie (5):

\[ \varphi_1\psi = \varepsilon_1\xi_1, \]

worin $\psi$ ungeändert geblieben ist. Folglich ist auch $\omega_1$ mit $\eta$ äquivalent, und mithin sind $\omega$ und $\omega_1$ unter einander äquivalent. % [sic] source prints $\eta$; repeating the derivation with unchanged $\psi$ requires $\eta'$
Das ist aber die in unserem Satze 3. ausgesprochene Behauptung, die also hiermit erwiesen ist.

Hiernach entspricht also jeder Idealclasse des Körpers $\Omega$ eine Classe äquivalenter Zahlen $\omega$ und umgekehrt, und wir können daher noch den Satz hinzufügen:
%CONT>>
